\documentclass[12pt]{article}

%%%%%%%%%%%%%%%%%%%%%%%%%%%%%%%%%%%%%%%%%%%%%%%%%%%%%%%%%%%%%%%%%%
%                     PREAMBLE / PACKAGES
%%%%%%%%%%%%%%%%%%%%%%%%%%%%%%%%%%%%%%%%%%%%%%%%%%%%%%%%%%%%%%%%%%

% ----------------------------------------------------------------
% Font encoding and page geometry
% ----------------------------------------------------------------
\usepackage[T1]{fontenc}
\usepackage[margin=1in]{geometry}

% ----------------------------------------------------------------
% Line and paragraph spacing
% ----------------------------------------------------------------
\usepackage{setspace}
\usepackage{parskip}

% ----------------------------------------------------------------
% Section and table-of-contents formatting
% ----------------------------------------------------------------
\usepackage{titlesec}
\usepackage{tocloft}

% ----------------------------------------------------------------
% Graphics, figures, and page layouts
% ----------------------------------------------------------------
\usepackage{graphicx}
\usepackage{animate} % frame-by-frame animations in supported PDF viewers
\usepackage[font=small]{caption}
\usepackage{subcaption}
\usepackage{float}
\usepackage{placeins}
\usepackage{wrapfig}
\usepackage{multicol}

% ----------------------------------------------------------------
% Mathematics
% ----------------------------------------------------------------
\usepackage{amsmath}
\usepackage{amssymb}
\usepackage{amsthm}
\usepackage{mathrsfs}
\usepackage{cancel}
\usepackage{gensymb}

% ----------------------------------------------------------------
% Tables and lists
% ----------------------------------------------------------------
\usepackage{booktabs}
\usepackage{tabularx}
\usepackage{multirow}
\usepackage{enumitem}

% ----------------------------------------------------------------
% Colors and highlighting
% ----------------------------------------------------------------
\usepackage{xcolor}
\usepackage{soul}
\sethlcolor{yellow}

% ----------------------------------------------------------------
% Headers and footers
% ----------------------------------------------------------------
\usepackage{fancyhdr}

% ----------------------------------------------------------------
% Source-code examples
% ----------------------------------------------------------------
\usepackage{listings}

% ----------------------------------------------------------------
% Citations and hyperlinks
% ----------------------------------------------------------------
\usepackage[numbers]{natbib}
\usepackage{hyperref}

% ----------------------------------------------------------------
% Default line spacing
% ----------------------------------------------------------------
\onehalfspacing
% ------------------------------------------------------------
% Document Information
% ------------------------------------------------------------
\title{Lab 6}
\author{Chuan Wei Heng}
\date{\today}

% ------------------------------------------------------------
% Document
% ------------------------------------------------------------
\begin{document}

\maketitle

\section*{Before lab}

We applied Newton's and Lazy Newton's method to the systems

\[
\mathbf{F}(x_1, x_2)
=
\begin{bmatrix}
x_1^2 + x_2^2 - 2 \\
e^{x_1-1} + x_2^2 - 2
\end{bmatrix}.
\]

We begin with the initial guesses $\mathbf{x}_0=(2,0.5)$
and $\mathbf{x}_0=(3,5)$. Interestingly, Newton's method
converged for the first initial guess but failed for the
second. As such, the Lazy Newton method was employed for the
second initial guess.

\begin{figure}[H]
    \centering

    \begin{subfigure}[t]{0.47\linewidth}
        \centering
        \includegraphics[width=\linewidth]{figures/pre_lab_newton.png}
        \caption{Newton iteration for $\mathbf{x_0}=(2, 0.5)$.}
        \label{fig:two_panel_a}
    \end{subfigure}
    \hfill
    \begin{subfigure}[t]{0.47\linewidth}
        \centering
        \includegraphics[width=\linewidth]{figures/pre_lab_lazy_newton.png}
        \caption{Lazy Newton iteration for $\mathbf{x_0}=(3, 5)$}
        \label{fig:two_panel_b}
    \end{subfigure}

    \caption{Comparison of convergence plots between Newton and Lazy Newton (Also Nikole, I took your suggestion and plotted classical convergence plots!)}
    \label{fig:two_panel}
\end{figure}

We can see that, while Newton's method failed to converge,
the Lazy Newton method did converge, although it required
significantly more iterations to do so.


\section*{Lab Day: Additional Quasi-Newton methods}

In the actual lab, we postulate conditions under which
the Jacobian should be recomputed. This modified method is
known as the ``Slacker Newton'' method. The system of interest
is given by

\[
\mathbf{F}(x_1, x_2)
=
\begin{bmatrix}
4x_1^2 + x_2^2 - 4 \\
x_1 + x_2 - sin(x_1-x_2)
\end{bmatrix}.
\]

We study the convergence at $\mathbf{x}0=(1,0)$ and
$\mathbf{x}0=(-1,0)$. The criterion chosen was the
absolute iteration error. We assign an arbitrary tolerance
denoted $tol{slacker}$ (for this lab, we set
$tol{slacker}=10^{-2}$) and recompute the Jacobian whenever
$||p_{n+1}-p_n||<tol_{slacker}$. The resulting convergence
behavior is shown below.

\begin{figure}[H]
    \centering

    \begin{subfigure}[t]{0.47\linewidth}
        \centering
        \includegraphics[width=\linewidth]{figures/slacker_newton_1.png}
        \caption{Roots = (0.999 -0.106)}
        \label{fig:two_panel_a}
    \end{subfigure}
    \hfill
    \begin{subfigure}[t]{0.47\linewidth}
        \centering
        \includegraphics[width=\linewidth]{figures/slacker_newton_2.png}
        \caption{Roots = (-0.999  0.106)}
        \label{fig:two_panel_b}
    \end{subfigure}

    \caption{Lazy Newton’s iteration with stopping criterion $\varepsilon=10^{-10}$.}
    \label{fig:two_panel}

\end{figure}

Unfortunately, I cannot compare results with my fellow
brethren from the lab, as I completed the write-up after the lab.
However, my friend told me that he used a relative error criterion
and achieved convergence in 4 iterations. This intuitively makes
sense, as in his code the Jacobian is likely recomputed more
frequently and, as such, fewer iterations may be required.

As for how this method performs on the example problem from class,
I genuinely do not know which example the prompt is referring to.
As such, I am defining an arbitrary 3D nonlinear system to test
the code on. The system is defined as follows:

\[
\mathbf{F}(x_1,x_2,x_3)
=
\begin{bmatrix}
4x_1^2+x_2^2-4\sqrt{x_3}\\
x_1x_3+x_2-\sin(x_1-x_2)\\
e^{x_3}+x_2^3+6x_1
\end{bmatrix}.
\]

\begin{figure}[H]

    \centering

    \includegraphics[width=0.75\linewidth]{figures/slacker_newton_for_random_system.png}

    \caption{Roots=(-1.056, 0.893, 1.73)}
    \label{fig:single_example}

\end{figure}

Overall, in this lab, we developed a modified Lazy Newton iteration, 
referred to as the Slacker Newton method. Rather than recomputing the 
Jacobian at every iteration, which can be computationally expensive, 
Slacker Newton updates the Jacobian only when necessary. This 
reduces the total number of Jacobian evaluations while retaining 
them at iterations where an update is most important.    

For details of the code, we turn the reader towards the git 
repository by \textit{chuanweiheng-cell} entitled APPM4600.
Thank you!

\end{document}